\subsection*{What We Found}
\begin{itemize}
  \item The original RF-DNA SVM pipeline contains three independent evaluation biases that collectively inflate reported ADR from $\svmOrigADR \to \svmADR$ (approximately 40 percentage points).
  \item Under a fair evaluation protocol, SVM and all three DL models
        (MLP-FV, GRU-DGT, BiGRU+TF-Aug) achieve statistically equivalent ADR in the 0.57--0.68 range. No architecture is a clear winner.
  \item With only five seeds the BiGRU shows ADR variance of $\pm 0.131$;
        LOTO cross-validation provides more stable per-device estimates
        (mean ADR \bigruLOTOadrMean, mean EER \bigruLOTOeerMean).
  \item Dataset size is the primary performance bottleneck: each binary
        classifier in trial\_1 trains on $\approx$52 windowed samples,
        far below what DL models typically require to generalise reliably.
\end{itemize}

\subsection*{Contributions}
\begin{itemize}
  \item \textbf{Revised baseline.} Identified, isolated, and quantified three evaluation assumptions in the original RF-DNA pipeline, providing reproducible corrections with per-bias ADR measurements.
  \item \textbf{Honest DL baseline.} First evaluation of MLP, GRU, and BiGRU architectures on this dataset under a strictly group-aware, transient-aware protocol across five seeds and LOTO.
  \item \textbf{Data-efficiency result.} Demonstrated that sliding-window augmentation enables DL training with 10$\times$ fewer recordings by exploiting the full transient duration, confirming that device-discriminating hardware imperfections persist well beyond the initial signal segment used by the original pipeline.
  \item \textbf{Multi-class probe.} Showed that a single BiGRU trained on all 12 devices reaches \mcAccuracy{} accuracy (vs.\ \mcRandomBaseline{} random), confirming that discriminative information exists across all devices and that architecture is not the limiting factor while dataset size is (Appendix~\ref{app:multiclass}).
\end{itemize}

\subsection*{Future Directions}
\begin{itemize}
  \item Collect additional independent transients per device to push ADR toward
        operationally useful levels.
  \item Evaluate AND-aggregation at system level to reduce the high False Accept
        Rate under OR-aggregation.
  \item Explore per-device classifier ensembles: training multiple models per
        authorized device and combining their outputs via soft voting could
        reduce the high seed-to-seed variance observed in this work and yield
        more stable authentication decisions without architectural changes.
  \item Conduct adversarial robustness evaluation (white-box and black-box
        attacks on raw IQ waveforms), now that an honest unattacked baseline
        is established.
\end{itemize}
